\chapter{社区参与式研究}
\chapterauthors{Alka M. Kanaya}
\sourcepage{485}
大多数临床研究在学术医疗中心或其他研究机构中构思和开展。这些场所具有许多优势，包括经验丰富的研究者和辅助人员。成熟的研究文化、声誉和基础设施，也有利于从初学者到终身教授的所有人开展工作。然而，这些便利可能伴随可推广性不佳、学术研究难以充分传播到社区等代价，甚至意外加重健康差异。为克服这些局限，美国通过国立卫生研究院的临床与转化科学奖励项目（Clinical Translational Science Awards，CTSA），以及以患者为中心的结局研究所（Patient-Centered Outcomes Research Institute，PCORI），大幅增加了支持社区参与研究的联邦经费。本章介绍背景各异的利益相关者如何合作，与社区共同设计、实施、解释、传播和应用研究成果。

我们把\term{社区参与式研究}定义为旨在满足开展研究所在社区需求的研究。“社区”可以有多种定义：居住在某一地理区域的人、有某种疾病的人、具有某种性别或族裔认同的人、非营利组织或倡导团体、同一卫生系统内的医疗服务者、政策制定者，或这些群体的组合。社区本位的研究方式要求相关社区利益相关者与研究中心的研究者合作。这种合作对于解决当地健康问题、促进健康公平至关重要，也能提供相互学习和建设能力的难得机会。
\sourcepage{486}
但它也充满挑战，因为建立信任、管理时间安排和预期、处理既有权力不平衡，都需要时间和精力。

\section{为什么开展社区参与式研究？}
涉及特殊场所、特定人群、新发或再发疾病的问题，合作研究往往是唯一的回答方式。学术医疗中心的研究重点，可能与当地社区需要相去甚远。参与研究给社区及学术研究者带来的益处，超出了某项研究所收集信息的价值。若以关切社区、促进健康公平的态度开展研究，便可能建立持久关系、产生自豪感，并发展经济和组织能力。

\subsection{当地问题与当地知识}
社区参与式研究可确保优先研究与当地密切相关的重要问题，如表15.1所示。来自中央渠道的国家或州级数据，未必反映当地疾病负担或社区危险因素分布。干预尤其是改变行为的干预，在不同场所未必效果相同。例如，戒烟公共卫生宣传可能无法有效覆盖越南裔美国人等吸烟率高的重要人口群体。因此，研究采用社区参与方法，先了解什么信息最有助于戒烟（1），再到最可能受益的社区检验这些针对文化特点的干预（2）。

\begin{bookbox}{表15.1\quad 需要社区参与式研究的问题示例}
\begin{itemize}
\item 芝加哥低收入社区青少年的电子烟使用率是多少？
\sourcepage{487}
\item 针对文化特点的外展服务能否提高亚裔美国人的结直肠癌筛查率？
\item 针对得克萨斯州流动农场工人的工作场所性传播感染预防宣传有何影响？
\item 社区卫生工作者干预能否改善波多黎各裔社区的2型糖尿病管理？
\end{itemize}
\end{bookbox}

当地知识可以识别新问题，找到更民主的解决办法（3）。不同社区对健康、医疗和疾病的理解与做法可能差异很大（4）。例如，造成健康差异的社会、政治和经济结构性障碍，可能使黑人男性难以在传统医疗场所获得最佳血压控制。一项学术机构与社区合作的研究，在洛杉矶县由黑人经营、备受信任的理发店中开展。研究让与顾客长期相识的理发师成为主要干预实施者（5）。在以理发店为随机化单位的整群随机试验中，血压控制不佳的顾客在理发店会见药师，由其开药并监测血压和电解质（干预组）；另一组由理发师鼓励改变生活方式并预约医生（积极对照组）。6个月后，干预组收缩压平均降幅比对照组大22 mm Hg，自评健康和患者参与指标也更好。

\subsection{提高可推广性}
社区研究有助于产生可推广到学术医疗中心代表不足人群的结果。例如，转诊医院的背痛患者与初级保健就诊者差异很大。因此，在三级医疗中心开展的背痛自然史或治疗反应研究，对社区临床实践的用途可能有限。
\sourcepage{488}
部分为应对这一问题，人们组建了多个\term{以实践为基础的研究网络}（practice-based research networks），让社区医生共同研究感兴趣的问题（6）。腕管综合征治疗反应就是一例。此前根据大型转诊中心患者的研究，建议尽早手术；但在初级保健中，大多数患者接受保守治疗即好转，很少需要转诊专科或复杂诊断检查（7）。

\subsection{建设当地能力与项目可持续性}
与社区利益相关者合作，可确保优先处理当地重要问题。社区参与的价值不止于具体信息，也会提高当地学术水平，鼓励创造力与独立思考。每个项目都能培养技能和信心，让当地研究者把自己视为科学过程的伙伴和领导者，而非仅仅消费外地知识。这又鼓励更多研究。此外，研究能为社区带来知识与经费，推动当地赋权和自立。

社区参与式研究还能实施和检验超越研究期限而持续的项目。让社区伙伴对干预内容、时间、方式和实施人员提出基础性意见，可增强社区的主人翁意识，获得更多认同，使项目长期延续。

\subsection{促进健康公平}
另一个重要理由是减少相关社区的\term{健康差异}（health disparities）（8）。合作中的多个环节——识别社区最关心的问题、寻找当地解决办法、建设有效且可持续项目的能力——都有助于促进\term{健康公平}（health equity）。
\sourcepage{489}
解决社区最迫切的问题，会增加信任和认同，改善结局、减少差异。最公正的方式，是社区、学术界及其他利益相关者共同收集和运用研究资料，立足社区优势及优先事项，改善健康与社会公平（9）。

社区研究还能改善卫生政策。例如，旧金山田德隆健康街角商店联盟开展集体行动，依据社区参与研究的发现，由社区组织、当地商家、学术伙伴、卫生部门和倡导团体共同努力，在该市最贫困的社区之一减少烟酒广告及销售，提高健康、可负担且可持续食品的可及性（10）。

\section{社区参与式研究的方法}
任何方法或设计都可以使用，从定性访谈、焦点小组和民族志观察，到流行病学调查或大型数据集的定量资料，再到临床试验。前述传统观察和实验设计，均可在研究全程吸收利益相关者意见；但也常需先评估社区需求来形成问题，并用开放式反馈帮助解释结果。多数项目采用混合方法，某些阶段需定性设计（第14章），另一些需定量设计。

理论上，启动项目与其他研究相同；实践中，年轻学术研究者最大的困难，是找到与社区已有联系、能够交流学习的同事或导师。当地或学术中心未必有这样的帮助。因此，初入这一领域者常需尽早决定：独立工作，还是与可能并无当地社区组织直接联系的资深研究者合作。
\sourcepage{490}
\subsection{如何起步}
没有经验丰富同事帮助而开始社区研究，就像自学游泳：并非不可能，却很困难，有时还有未预料的危险。以下几步可使起步容易一些。
\begin{itemize}
\item \textbf{建立联系网络。}如第2章所述，这对每位研究者都重要。通常，与已在目标社区建立信任关系的资深研究者合作，是良好起点。新研究者应联系研究类似问题或与目标社区合作的人。参加领域会议有助于建立联系、发现当地志趣相投的同事。许多学术中心的临床研究设施包含社区参与办公室，工作人员可帮助新研究者联系当地和区域伙伴。寻找研究过相关主题的团体或项目，也能通向其他社区资源。
\item \textbf{从简单项目开始。}初到社区就开展随机对照试验，通常并非好主意。能提供有用当地资料的小型描述性研究往往更适合，也能为日后合作奠定基础。雄心更大的项目可以留待以后。例如，芝加哥南亚移民的心脏病知识描述性研究（11），成为该社区更大型心血管病预防行为干预的第一步（12）。
\item \textbf{考虑当地优势。}当地有什么问题比别处更适合回答？年轻研究者通常应聚焦社区常见的健康问题或人群，并征求其意见，确定最相关、最重要的问题。
\end{itemize}

\sourcepage{491}
\subsection{合作研究的连续谱}
应考虑社区参与的强度（图15.1）（13）。
\begin{figure}[htbp]
\centering
\begin{tikzpicture}
\filldraw[draw=dcrblue!60,fill=dcrblue!45] (0,0) -- (14.6,0) -- (14.6,-.2) -- (15.3,.4) -- (14.6,1) -- (14.6,.8) -- (0,.8) -- cycle;
\node[font=\bfseries] at (7.3,.4) {社区参与强度};
\end{tikzpicture}
\par\vspace{3pt}
\begin{minipage}[t]{.31\textwidth}\small\centering\textbf{有限参与}\par\raggedright
\begin{itemize}[leftmargin=1.1em]
\item 协助具体任务
\item 招募
\item 社区外展
\item 社区参与工作坊
\end{itemize}\end{minipage}\hfill
\begin{minipage}[t]{.31\textwidth}\small\centering\textbf{实质性参与}\par\raggedright
\begin{itemize}[leftmargin=1.1em]
\item 为设计和方法提供意见
\item 招募
\item 在社区实施研究
\item 社区咨询委员会
\end{itemize}\end{minipage}\hfill
\begin{minipage}[t]{.31\textwidth}\small\centering\textbf{全面参与}\par\raggedright
\begin{itemize}[leftmargin=1.1em]
\item 全面平等的伙伴关系
\item 确定问题、撰写申请、争取经费
\item 共同开展研究设计、方法制定、实施、分析和传播
\end{itemize}\end{minipage}
\caption{社区参与式研究的连续谱。}
\end{figure}

有限参与的研究由学术中心发起，社区伙伴协助招募。例如，研究者与社区伙伴及工作人员共同开发适合文化和语言的招募资料。在这一层次，学术研究者通常负责设计、争取经费和取得实施许可。可通过社区参与工作坊，在规划早期征求意见（表15.2）。社区伙伴获得支持外展和招募的资金与资源，并积累研究经验。

\begingroup\small
\begin{longtable}{L{.23\textwidth}L{.70\textwidth}}
\caption*{\textbf{表15.2\quad 社区参与式研究的具体方式}}\\\toprule
方式 & 内容\\\midrule\endfirsthead\toprule 方式 & 内容\\\midrule\endhead
社区外展 & 研究团队的社区及学术成员开展讲座，提高公众对社区重要问题的认识。可以吸引社区成员参与，帮助招募和留存，产生未来研究的新想法。\\
\sourcepage{492}社区参与工作坊 & 一种结构化过程，可由机构提供服务，便于相关社区和患者利益相关者针对具体项目提出意见，改善设计、实施和传播（14）。研究者向社区成员组成的小组作简短介绍，并提出几个问题；成员反馈，实现共同学习和过程评价。\\
社区咨询委员会（community advisory boards，CAB） & 充当社区与研究者之间的联络者。通常包括不同社区群体代表，如社区组织服务对象和工作人员、宗教团体、学校及媒体成员。研究早期即组建，成员审阅研究各部分、提供意见，并在整个研究期间定期会面。\\\bottomrule
\end{longtable}\endgroup

实质性参与的研究让利益相关者共同设计和实施方案。社区伙伴常作为经费分包方，与学术团队合作，在社区实施教育或咨询等干预。研究全程常设社区咨询委员会，对方案、招募资料等提供意见；委员会还参与传播结果及规划下一步合作（表15.2）。这一程度参与的优势包括：在社区研究本身的科学价值、成果共同署名、建立合作伙伴关系的满足感，以及帮助社区发展研究能力。

Live Well Be Well（“好好生活，保持健康”）研究（15）是中等强度参与的例子，由学术机构与社区组织共同治理、分工负责。社区伙伴是伯克利市加州公共卫生部门（Department of Public Health，DPH）的工作人员，参与糖尿病预防试验方法设计，实施电话干预，协助传播发现。学术伙伴为加利福尼亚大学旧金山分校研究者，协助外展和招募，负责项目评价及数据分析。
\sourcepage{493}
该随机对照试验发现，电话生活方式干预使社区中糖尿病前期参与者的体重和甘油三酯显著下降，饮食也得到改善（16）。由于卫生部门从最初试验阶段就参与，项目在建立时即考虑了持续运行和补充其常规服务，试验结束多年后仍在开展。

全面参与指学术研究者与利益相关者合作建设研究基础设施，寻找互惠机会。在最为周密的模式中，社区参与从构思到传播发现的所有环节。\term{社区本位参与式研究}（community-based participatory research，CBPR）以社区优先事项为基础，由社区与学术研究者共同决策（9）。最终目标和回报是改善社区结局与健康公平。长期 CBPR 项目包括：旧金山亚裔美国人癌症意识、研究与培训网络，合作促进癌症认识和预防（17）；北卡罗来纳州农村非裔美国人社区的 GRACE 项目，起初关注 HIV 认识和预防，后来扩展到心血管病预防（18，19）。

许多 PCORI 资助研究明确要求患者、家庭、社区成员和医疗利益相关者参与项目开发、治理、实施和传播的全部环节（20）。例如，在超重且患高血压或2型糖尿病者中，一项三组整群随机试验由患者和医疗利益相关者（负责初级保健诊所常规患者联系的非临床人群健康管理人员）共同选择、完善减重干预。试验比较常规照护、单独在线体重管理，以及在线项目加人群健康管理人员的联合干预。
\sourcepage{494}
联合干预的减重效果优于在线项目或常规照护（21）。该试验的优势在于利用已融入诊所常规照护的人群健康团队，并征求工作人员意见，使项目嵌入其工作流程。试验表明，初级保健诊所的系统层面变革可行且有效，能够扩大并持续实施，以改善患者健康。

\section{社区参与式研究的挑战}
如果不深入理解社区的文化视角，即使计划周密、技术先进，研究也可能失败。研究者必须理解当地对疾病的认识和做法，发展在文化上适当的合作方式。沟通不频繁、询问迟迟得不到答复，都是合作可能出问题的征兆。

\subsection{建立信任}
学术与社区伙伴之间的信任是基础，而且应是双向的。\textbf{双方必须相互信任，都相信合作能够互惠。}建立研究议程前，就应先建立伙伴关系。个人、人际、社会和结构各层面的因素都可能损害信任。许多社区过去曾遭学术研究者忽视、居高临下对待或利用，留下不信任的历史（22）。学术研究者应有意识地承认过错、坦诚对话、保持谦逊，尊重所有声音和经历，修复并加强关系（23）。

学术研究者通常处于有权力、有特权的位置：来自资源丰富的机构，具备临床研究经验、训练和其他实际支持。
\sourcepage{495}
承认这些优势，并与社区伙伴分享知识、技能、机会和经费，是长期信任的必要条件。

建立信任需要时间与努力。与已同社区建立信任的同事合作会有帮助；但年轻研究者也应建立自己独立、清晰、一贯的沟通渠道，显示对伙伴的承诺。定期安排会议，让双方交流进展、提出困难和下一步意见，是一种办法；最好备有食物饮料，并留出社交时间。

\subsection{管理时间安排和预期}
与社区伙伴、利益相关者安排规划会议可能很困难，常需在周末或晚间举行。项目开始后，社区伙伴还需在组织的整体活动中腾出研究所需的时间与空间。这些事应在规划早期商定，否则可能打乱内部或资助者要求的时间表。一般应预期启动和招募会延误，并及早、持续地征询社区利益相关者，根据其经验判断最可行的安排。

学术界和多数社区都有权力等级，但有些社区没有正式组织结构，难以寻找伙伴。此时应寻求与社区接触的倡导团体或社会服务机构。不过，如果代表不能体现最弱势者的利益，合作可能加重不平等，而非增强社区能力。

新伙伴关系还应明确各方角色与预期。多数以学术机构为主的临床研究，会为各分包场所规定工作范围、任务和交付内容。
\sourcepage{496}
尤其没有正式分包合同时，必须清楚交流预期。此时可考虑共同制定、签署\term{谅解备忘录}（Memorandum of Understanding，MOU），记录工作范围、交付内容及大致时间表。

学术界通过经费，以及论文和报告的数量、质量来奖励科学产出；非学术伙伴的回报则可能是可持续的健康改善项目、就业机会和差异缩小。研究者若使自己的优先事项、长期目标与社区伙伴相契合，能获得最多回报和个人成长；学术医疗中心若也使内部文化更灵活、更包容社区意见，效果尤其明显（24）。加强合作、拓展研究机会，能让所有利益相关者受益。

\section{小结}
\begin{enumerate}
\item 社区参与式研究可识别文化及其他当地因素，判断特定社区最需要、且可能有效的干预。
\item 参与临床研究能产生更适用的发现，建设当地能力和可持续项目，促进健康公平，使社区伙伴受益。
\item 设计和伦理问题虽与其他临床研究相似，建立伙伴关系、寻找导师等实际问题在社区可能更难。应建立联系网络，从小项目开始，发现当地优势。
\item 学术中心与社区研究者的合作强度呈连续谱：有限参与时，社区协助学术伙伴发起的具体任务；实质性参与时，社区协作开展外展和实施；全面参与时，双方平等合作研究的全部环节。
\sourcepage{497}
\item 具体方式包括社区外展、社区参与工作坊和社区咨询委员会。
\item 挑战包括建立信任所需的时间，以及妥善管理时间安排和预期。
\end{enumerate}

\begin{references}
\item Kenny JD, Tsoh JY, Nguyen BH, Le K, Burke NJ. Keeping each other accountable: social strategies for smoking cessation and healthy living in Vietnamese American men. Fam Community Health. 2021;44(3):215-224.
\item Tong EK, Saw A, Fung LC, Li CS, Liu Y, Tsoh JY. Impact of a smoke-free-living educational intervention for smokers and household nonsmokers: a randomized trial of Chinese American pairs. Cancer. 2018;124(Suppl 7):1590-1598.
\item Corburn J. Bringing local knowledge into environmental decision making: improving urban planning for communities at risk. J Plan Educ Res. 2003;22(4):420-433.
\item Griffith BN, Lovett GD, Pyle DN, et al. Self-rated health in rural Appalachia: health perceptions are incongruent with health status and health behaviors. BMC Public Health. 2011;11:229.
\item Victor RG, Lynch K, Li N, et al. A cluster-randomized trial of blood-pressure reduction in black barbershops. N Engl J Med. 2018;378:1291-1301.
\item Nutting PA, Beasley JW, Werner JJ. Practice-based research networks answer primary care questions. JAMA. 1999;281:686-688.
\item Miller RS, Ivenson DC, Fried RA, et al. Carpal tunnel syndrome in primary care: a report from ASPN. J Fam Pract. 1994;38:337-344.
\item Cooper L. Rethink how we plan research to shrink COVID health disparities. Nature. 2021;590:9.
\item Wallerstein N, Duran B, Oetzel J, Minkler M, eds. Community-Based Participatory Research for Health: Advancing Social and Health Equity. 3rd ed. Jossey-Bass; 2018.
\item Flood J, Minkler M, Hennessey Lavery S, et al. The Collective Impact Model and its potential for health promotion: overview and case study of a healthy retail initiative in San Francisco. Health Educ Behavior. 2015;42(5):654-668.
\item Kandula NR, Tirodkar MA, Lauderdale DS, et al. Knowledge gaps and misconceptions about coronary heart disease among U.S. South Asians. Am J Prev Med. 2010;38(4):439-442.
\sourcepage{498}
\item Kandula NR, Patel Y, Dave S, et al. The South Asian Heart Lifestyle Intervention (SAHELI) study to improve cardiovascular risk factors in a community setting: design and methods. Contemp Clin Trials. 2013;36(2):479-487.
\item Pasick R, Oliva G, Goldstein E, Nguyen T. Community-engaged research with community-based organizations: a resource manual for UCSF researchers. In: Fleisher P, ed. UCSF Clinical and Translational Science Institute (CTSI) Resource Manuals and Guides to Community-Engaged Research. Clinical Translational Science Institute Community Engagement Program, University of California San Francisco; 2010. http://ctsi.ucsf.edu
\item Joosten YA, Israel TL, Williams NA, et al. Community Engagement Studios: a structured approach to obtaining meaningful input from stakeholders to inform research. Acad Med. 2015;90(12):1646-1650.
\item Delgadillo AT, Grossman M, Santoyo-Olsson J, et al. Description of an academic-community partnership lifestyle program for lower-income, minority adults at risk for diabetes. Diabetes Educ. 2010;36(4):640-650.
\item Kanaya AM, Santoyo-Olsson J, Gregorich S, et al. A telephone-based lifestyle intervention trial to lower risk factors in ethnic minority and lower socioeconomic status adults at risk of diabetes: the live well be well study. A randomized controlled trial. Amer J Public Health. 2012;102(8):1551-1558.
\item McPhee SJ, Nguyen TT, Mock J, et al. Highlights/best practices of San Francisco’s Asian American Network of Cancer Awareness, Research, and Training (AANCART). Cancer. 2005;104(12):2920-2925.
\item Corbie-Smith G, Adimore AA, Youmans S, et al. Project GRACE: a staged approach to development of a community-academic partnership to address HIV in rural African American communities. Health Promot Pract. 2011;12(2):293-302.
\item Corbie-Smith G, Wiley-Cene C, Bess K, et al. Heart Matters: a study protocol for a community-based randomized trial aimed at reducing cardiovascular risk in a rural, African American community. BMC Public Health. 2018;938.
\item Hickam D, Totten A, Berg A, et al. The PCORI Methodology Report. Patient-Centered Outcomes Research Institute; 2013.
\item Baer HJ, Rozenflum R, De La Cruz BA, et al. Effect of an online weight management program integrated with population health management on weight change. A randomized clinical trial. JAMA. 2020;324(17):1737-1746.
\item Pacheco CM, Daley SM, Brown T, et al. Moving forward: breaking the cycle of mistrust between American Indians and researchers. Am J Public Health. 2013;103:2152-2159.
\sourcepage{499}
\item Moreno-John G, Gachie A, Fleming CM, et al. Ethnic minority older adults participating in clinical research: developing trust. J Aging Health. 2004;16(5 Suppl):93S-123S.
\item Michener L, Cook J, Ahmed SM, et al. Aligning the goals of community-engaged research: why and how academic health centers can successfully engage with communities to improve health. Acad Med. 2012;87(3):285-291.
\end{references}

\chapterappendix{15A}{第15章习题：社区参与式研究}
\sourcepage{500}
\begin{enumerate}
\item 一名新研究者希望改善亚裔移民女性的乳腺癌筛查，目前处于早期规划阶段，正考虑如何让社区伙伴参与。
\begin{enumerate}
\item 描述可用于有限或实质性社区参与的不同方法。
\item 两种参与强度各有哪些优势和挑战？
\end{enumerate}
\end{enumerate}
